\section{Priority Classification}

\subsection{Classification Rules}

The P1/P2/P3 classification uses two signals: \textbf{reviewer count} (how many reviewers raised the issue) and \textbf{severity} (major vs.\ minor in the original review).

\begin{table}[h]
\centering
\begin{tabular}{lcl}
\toprule
Priority & Rule & Rationale \\
\midrule
\textbf{P1} & 3+ reviewers OR any marks ``major'' & Broad agreement or blocking severity \\
\textbf{P2} & 2 reviewers & Meaningful but not universal concern \\
\textbf{P3} & 1 reviewer, minor & Individual preference \\
\bottomrule
\end{tabular}
\caption{Priority classification rules.}
\end{table}

\subsection{Rule Justification}

The 3-reviewer threshold for P1 is calibrated to panels of 5--9 reviewers, where 3+ represents a majority or near-majority. The ``any major'' override ensures that a single reviewer's blocking concern is not downgraded by the lack of other reviewers in that specific area.

\subsection{Threshold Sensitivity Analysis}

To verify that the 3-reviewer threshold is robust rather than an arbitrary choice, we sweep the reviewer-count threshold from 1 to $N{-}1$ and measure the resulting P1 classification's score impact across all 33 synthesis cycles.

\begin{table}[h]
\centering
\small
\begin{tabular}{lcccc}
\toprule
Threshold & P1 Items/Paper & Score Impact & \% of Total Improvement & Precision$^\dagger$ \\
\midrule
$\geq 1$ (any mention) & 8.7 & +1.8 & 94\% & 0.41 \\
$\geq 2$ & 5.4 & +1.5 & 83\% & 0.67 \\
\textbf{$\geq 3$ (selected)} & \textbf{2.8} & \textbf{+1.3} & \textbf{72\%} & \textbf{0.91} \\
$\geq 4$ & 1.2 & +0.7 & 39\% & 0.96 \\
$\geq 5$ & 0.4 & +0.3 & 17\% & 0.98 \\
\bottomrule
\end{tabular}
\caption{P1 classification sensitivity to reviewer-count threshold. $^\dagger$Precision measured against author-annotated ground truth P1 items (87\% agreement rate). The selected threshold ($\geq 3$) achieves the best balance between coverage (72\% of improvement) and precision (0.91).}
\label{tab:threshold-sensitivity}
\end{table}

The threshold of 3 sits at the ``elbow'' of the precision-coverage curve: lowering it to 2 increases the P1 count by 93\% while gaining only 11 percentage points of coverage, flooding authors with low-priority items. Raising it to 4 maintains high precision but misses a third of the score improvement. This pattern is stable across panel sizes: for panels of 5 ($N{=}5$, threshold $\geq 3$ captures 60\% of reviewers), 7 ($\geq 3$ captures 43\%), and 9 ($\geq 3$ captures 33\%). Larger panels may benefit from a higher absolute threshold, though the relative position ($\sim$50\% of panel) remains effective. We discuss adaptive thresholds in Section~\ref{sec:discussion}.

\subsection{Distribution Analysis}

Across 33+ synthesis cycles:

\begin{table}[h]
\centering
\begin{tabular}{lccc}
\toprule
Priority & Avg Items/Paper & Avg Reviewer Count & \% of Total Issues \\
\midrule
P1 & 2.8 & 3.4 & 23\% \\
P2 & 4.2 & 2.1 & 35\% \\
P3 & 5.1 & 1.0 & 42\% \\
\bottomrule
\end{tabular}
\caption{Priority distribution across 33+ synthesis cycles.}
\end{table}

The pyramid structure (few P1, many P3) is desirable: most issues are minor, and only a small number require immediate attention. This focuses author effort on the vital few rather than the trivial many.

\subsection{Common P1 Issue Types}

The most frequent P1 issues across all papers:

\begin{enumerate}
    \item \textbf{Missing competitive comparison} (79\% of papers)
    \item \textbf{Generalization concerns} (64\% of papers)
    \item \textbf{Statistical rigor gaps} (50\% of papers)
    \item \textbf{Overclaimed contributions} (36\% of papers)
    \item \textbf{Reproducibility deficiencies} (29\% of papers)
\end{enumerate}
